\begin{textsample}{POS Dim 1 – sc – Score 75.00 – sc/sc\_em\_60.txt}  \label{ex:f1_pos_008}
3.2 \textbf{TRILHA} 1 : “ DIÁLOGOS COM NOSSAS CIDADES - MEIO AMBIENTE E SUSTENTABILIDADE ” Área do conhecimento : Ciências da Natureza e suas \textbf{tecnologias}.
Carga horária : 160 h ou 240 h, a \textbf{depender} da \textbf{matriz} curricular em funcionamento na Unidade Escolar. \textbf{Aulas} semanais : 10 ou 15 \textbf{aulas}, conforme \textbf{matriz} da escola.
Perfil docente : Três professores licenciados, ao menos um para cada componente da Área de Ciências da Natureza e suas \textbf{tecnologias}, que trabalhem de \textbf{maneira} \textbf{articulada} : - Habilitação em Ciências Biológicas : 3 \textbf{aulas} ou 5 \textbf{aulas}, a \textbf{depender} da \textbf{matriz} curricular em funcionamento na Unidade Escolar ; - Habilitação em Física : 4 \textbf{aulas} ou 5 \textbf{aulas}, a \textbf{depender} da \textbf{matriz} curricular em funcionamento na Unidade Escolar ; - Habilitação em Química : 3 \textbf{aulas} ou 5 \textbf{aulas}, a \textbf{depender} da \textbf{matriz} curricular em funcionamento na Unidade Escolar.
A \textbf{trilha} “ Diálogos com nossas cidades - meio ambiente e sustentabilidade ” é um convite para conhecer e pensar quem somos como cidadãos, o que estamos fazendo e o que podemos fazer para deixar nossa cidade melhor.
Ao olhar para as cidades, \textbf{inúmeros} questionamentos podem ser feitos.
Entre eles : - Como é o lugar em que você vive e sua qualidade de vida? - Estaria você desfrutando de um alto padrão de qualidade de vida, vivendo em um lugar com planejamento urbano, serviços públicos de qualidade, acesso à \textbf{tecnologia}, ou estaria você enfrentando \textbf{problemas} como pobreza, criminalidade e desemprego? - Para onde e para que estamos seguindo com um desenvolvimento acelerado ( na maioria das vezes, sem planejamento )? - O que seria de uma cidade sem energia elétrica? - O que são os veículos e que combustíveis os movimentam? - O que havia nos lugares em que construímos as cidades de que nos orgulhamos?
Estes são alguns dos questionamentos que podem ser respondidos com o criterioso olhar das Ciências da Natureza e suas \textbf{tecnologias} e que motivaram a criação desta \textbf{trilha}.
A \textbf{compreensão} dos objetos de conhecimento contemplados nas unidades curriculares proporciona a estudantes e professores um novo olhar sobre a dinâmica das cidades, com possibilidades de propor alternativas para \textbf{problemas} locais e a realização de interferência direta na comunidade.
Ao olhar pela janela de nossas cidades, usando as lentes das Ciências da Natureza e suas \textbf{tecnologias}, o campo de visão se amplia e novos \textbf{horizontes} se abrem. 3.2.1 Justificativa As primeiras cidades desenvolveram -se na Mesopotâmia, em torno de 8.000 a 3.500 anos a.C.
Tiveram origem a partir de aglomerados mais densamente habitados, \textbf{influenciados} pelo comércio e centros militares.
Não é o propósito da \textbf{trilha} explorar as transformações que as cidades vieram sofrendo com o passar de todos estes anos, mas promover uma reflexão sobre como era a região onde o sujeito nasceu e como ela se encontra agora.
Temos que admitir que muitas cidades estão em ritmo de crescimento. \textbf{Inúmeras} são as transformações em diversas perspectivas que acontecem num curto período de tempo.
Por um \textbf{lado}, a população humana \textbf{segue} aumentando, assim como o desejo pelo lucro, colocado, geralmente, como objetivo \textbf{primordial}.
Por outro \textbf{lado}, há reflexos no desenvolvimento das cidades que afetam as pessoas que nelas vivem, e que podemos analisar.
Pietrocola ( 2019 ) destaca que nos últimos 75 anos a \textbf{humanidade} começou a se \textbf{preocupar} mais com o que fizemos com a natureza, deixando claro que as cidades não são rodeadas por natureza, mas que a natureza foi cercada pela \textbf{humanidade}.
A sociedade pós-natureza reflete e experimenta como as \textbf{ciências} e a \textbf{tecnologia} transformaram a natureza em “ natureza tecnológica ”, \textbf{tanto} em sentido positivo, quanto no sentido de analisar os riscos ( entre eles o da degradação ambiental, do aumento da pobreza, da ampliação das desigualdades ) provocados pelas transformações.
Nossas cidades enfrentam \textbf{inúmeros} \textbf{problemas} relacionados ao congestionamento do tráfego, à infraestrutura urbana inadequada e à ausência de serviços básicos, como saneamento básico e abastecimento de água ( ou à escassez dela também ), à gestão de resíduos sólidos e ao controle da poluição.
Dessa forma, faz -se necessário repensar nossas relações com o ambiente e com as cidades, objetivando a construção de uma sociedade cada vez mais sustentável, em que a educação \textbf{desempenhe} papel fundamental na \textbf{tarefa} de formar cidadãos \textbf{críticos} e aptos a atuar na realidade em que estão inseridos e a melhorá -la.
Cidades sustentáveis são capazes de propiciar um padrão de vida aceitável, sem causar profundos prejuízos ao ecossistema ou aos ciclos biogeoquímicos de que ela \textbf{depende} ( MARK ROSELAND, 1997 ).
O conjunto de unidades curriculares da \textbf{trilha} de \textbf{aprofundamento} “ Diálogos com nossas cidades – Meio Ambiente e Sustentabilidade ” poderá \textbf{auxiliar} na construção dos projetos de vida dos(as) estudantes, uma vez que possibilita conhecer \textbf{caminhos} em relação a questões sociais, ambientais, de saúde pública e do mundo do trabalho, e sobre elas refletir e investigar, além de propor alternativas, almejando a melhoria da qualidade de vida e a construção de uma sociedade cada vez melhor.
Neste sentido, busca -se \textbf{aproximar} a complexidade das Ciências da Natureza e suas \textbf{tecnologias} da vida dos(as) estudantes, numa proposta \textbf{dialógica} e engajada, para que possam formar suas opiniões e gerenciar os riscos de suas decisões, como defende Pietrocola ( 2019 ).
Assim, a presente \textbf{trilha} poderá contribuir de \textbf{maneira} significativa para o desenvolvimento local e regional, pois possibilita que os(as) estudantes, a partir do conhecimento aprofundado de aspectos da realidade \textbf{socioambiental} de sua cidade e região, de suas diversidades ( sociais, culturais, \textbf{econômicas} e ambientais ), possam exercer seu protagonismo juvenil.
O(a) estudante poderá adquirir conhecimentos de Biologia, Física e Química para sustentar a investigação dos principais \textbf{problemas} que o cercam, elaborando e testando hipóteses, buscando parcerias e promovendo discussões junto às empresas, ao poder público e a outras entidades acerca das questões sociais e ambientais para o desenvolvimento socioeconômico local e regional.
Podem ser consideradas as particularidades e potencialidades da cidade em que se localiza a escola, porém, de \textbf{maneira} cada vez mais sustentável, visando à melhoria da qualidade de vida, \textbf{tanto} em nível individual quanto coletivo, a curto e a longo prazo.
Entre os aspectos que serão aprofundados nessa \textbf{trilha} estão : recursos energéticos locais ; formas de produção, distribuição e consumo de energia elétrica ; utilização racional e sustentável dos recursos naturais ; tratamento de água e esgoto ; gerenciamento de resíduos sólidos produzidos ; controle da erosão urbana ; mobilidade urbana ; \textbf{problemas} de saúde ocasionados pela dinâmica das cidades ; consumo consciente ; princípios da sustentabilidade e plano diretor municipal.
Propõe -se que a \textbf{trilha} seja desenvolvida utilizando todos os recursos tecnológicos possíveis, através de metodologias ativas que permitam que o estudante participe ativamente do processo de ensino e aprendizagem e se aproprie do conhecimento.
Além disso, considerando que cada município apresenta características próprias, com diferentes públicos, diferentes questões culturais, \textbf{econômicas} e ambientais, os objetos de conhecimento propostos em cada unidade curricular devem ser trabalhados considerando e respeitando a realidade de cada estudante, de cada escola e de cada comunidade. 3.2.1 Objetivo Geral Utilizar conhecimentos de Biologia, Física e Química para compreender a dinâmica das cidades por meio da investigação de \textbf{problemas} ligados ao planejamento urbano, à eletricidade, à poluição, à mobilidade urbana, aos recursos naturais para que os(as) estudantes conheçam as cidades e proponham alternativas para a melhoria das condições ambientais e para a solução de \textbf{problemas} locais. 3.2.3 Desenvolvimento A \textbf{trilha} de \textbf{aprofundamento} “ Diálogos com as cidades - Meio Ambiente e Sustentabilidade ” aborda o conhecimento científico associado às diferentes relações entre os seres humanos, com os demais seres vivos e com os ambientes em que vivem, na busca de melhores condições de sobrevivência, tendo como \textbf{foco} as cidades.
Propõe experiências, técnicas, conhecimentos e valores que, produzidos na e pela coletividade, possam interferir sobre a realidade vivida, com implicações nos fatores \textbf{socioambientais}.
Apresenta, em seu \textbf{escopo}, unidades curriculares que conversam entre si, trazendo vários enfoques sobre temas \textbf{contemporâneos} e extremamente importantes para a vida das cidades, com uma organização flexível, permitindo à escola optar pela \textbf{sequência} de trabalho das unidades.
Nas orientações \textbf{metodológicas}, aponta possibilidades de trabalho, sendo fundamental que os docentes reflitam sobre elas e definam sua viabilidade, de acordo com suas realidades locais, podendo produzir novos designs \textbf{metodológicos} que se ajustem melhor ao processo de ensino e aprendizagem.
A organização em eixos estruturantes traz conhecimentos científicos à área de \textbf{Ciências} da Natureza e suas \textbf{tecnologias}, considerando o mundo \textbf{contemporâneo} em que uma chamada sociedade da informação, com cenários de rapidez tecnológica, \textbf{exige} atuação consciente e responsável dos(as) estudantes.
Assim, as habilidades apresentadas \textbf{supõem} o desenvolvimento de processos de ensino e aprendizagem que favoreçam o protagonismo, incentivem a pesquisa e valorizem o empreendedorismo na realização de projetos de vida.
O \textbf{aprofundamento} dos conhecimentos científicos propostos na \textbf{trilha} está em consonância com o proposto por Moran : > [... ] Ensinar e aprender tornam -se fascinantes quando se convertem em processos de pesquisa constantes, de questionamento, de criação, de experimentação, de reflexão e de compartilhamentos crescentes, em áreas de conhecimento mais amplas e em níveis cada vez mais profundos [ … ] ( 2018, p. 3 ).
Isto porque a proposta possibilita a investigação da realidade local para ir além do contexto vivido e reconhecer -se como partícipe de uma sociedade, responsável pelo ambiente e por sua sobrevivência.
Esta \textbf{trilha} foi organizada em quatro unidades curriculares, conforme \textbf{figura} 3. \textbf{Figura} 3 - Unidades curriculares Fonte : Elaboração dos \textbf{autores} ( 2020 ).
Unidade Curricular I : Energia elétrica - da produção ao consumo Objetivo : Promover um processo \textbf{reflexivo} acerca das manifestações e transformações de energia, produção, distribuição e consumo de energia elétrica, de sua importância para a sociedade e dos impactos ambientais decorrentes de sua produção e utilização, incentivando o protagonismo juvenil na busca por alternativas de produção e consumo mais sustentáveis.
Carga horária : total de 40 h ou 60 h, a \textbf{depender} da \textbf{matriz} curricular em funcionamento na Unidade Escolar. \textbf{Quadro} 26 - Energia elétrica - da produção ao consumo Eixo estruturante/Palavras-chave das habilidades específicas - Processos criativos - Refletir criativa/criticamente - Energia - Selecionar e \textbf{mobilizar} conhecimentos científicos Objetos do conhecimento - Manifestações de energia : cinética, \textbf{potencial}, química, térmica, sonora, luminosa, elétrica. - Discussão sobre as fontes energéticas e os impactos ambientais e sociais da geração e utilização da energia nos diferentes setores da sociedade. - Química verde : possibilidades de redução do uso de fontes não renováveis de energia em nível local e global. - Bioluminescência : produção de luz por seres vivos ( fungos, insetos, águas-vivas, peixes, algas, entre outros ), conceitos, importância ecológica e aplicações. - Energia elétrica - Produção de energia elétrica ( em nível nacional e local ) a partir de fontes renováveis e não renováveis de energia : eficiência energética e impactos ambientais. - Diferenciação entre corrente alternada e corrente contínua, e sua relação com a geração de energia elétrica em grande escala. - Leis e processos envolvidos na produção ( geradores ), distribuição ( \textbf{transformadores} ) e consumo ( motores ) de energia elétrica. - Distribuição de energia elétrica e utilização nos diferentes setores da sociedade : percentual de energia utilizado nas indústrias, no comércio e nas residências. - Aparelhos elétricos e consumo de energia elétrica ( transformações de energia ; grandezas elétricas que a caracterizam como corrente elétrica, resistência, tensão e potência ). - Consumo e custo da energia elétrica ( leitura e interpretação de uma conta de energia elétrica - “ conta de luz ” - consumo, tarifas e taxas ). - \textbf{Influência} da sociedade de consumo na produção de energia elétrica e no desenvolvimento de equipamentos elétricos. - Impactos negativos da luminosidade no céu noturno das cidades, e seus impactos na fauna local.
Habilidades de \textbf{aprofundamento} da área - Analisar e representar, com ou sem o uso de dispositivos e de \textbf{aplicativos} \textbf{digitais} específicos, as transformações e conservações em sistemas que envolvam quantidade de matéria, de energia e de movimento para realizar previsões sobre seus comportamentos em situações cotidianas e em processos produtivos que priorizem as sociedades sustentáveis, o uso consciente dos recursos naturais e a preservação da vida em todas as suas formas. - Analisar questões \textbf{socioambientais}, políticas e \textbf{econômicas} relativas à dependência do mundo \textbf{atual} em relação aos recursos não renováveis e discutir a necessidade de se introduzir alternativas e novas \textbf{tecnologias} energéticas e materiais, comparando diferentes tipos de motores e processos de produção de novos materiais. - Estimar o consumo de energia elétrica de uma residência e propor e \textbf{divulgar} alternativas para uma utilização sustentável. - Pesquisar e elaborar produtos de divulgação científica ( \textbf{simulação}, texto, \textbf{vídeo}, reportagem ) sobre produção, distribuição, impactos ambientais e/ou consumo de energia elétrica na cidade.
Fonte : Elaboração dos \textbf{autores} ( 2020 ).
Unidade Curricular II : Impactos \textbf{socioambientais} locais e globais relacionados ao solo, ao ar e à água Objetivo : Refletir acerca dos \textbf{problemas} e impactos \textbf{socioambientais} relacionados ao solo, ao ar e à água em nível local e global, suas principais causas e consequências, de forma a investigar a realidade local por meio do desenvolvimento de projetos de pesquisa elaborados coletivamente na busca por alternativas para diminuir os impactos sociais e ambientais locais decorrentes da utilização inadequada dos recursos naturais.
Carga horária : Total de 40 h ou 60 h, a \textbf{depender} da \textbf{matriz} curricular em funcionamento na Unidade Escolar. \textbf{Quadro} 27 - Impactos \textbf{socioambientais} locais e globais relacionados ao solo, ao ar e à água Eixo estruturante/Palavras-chave das habilidades específicas - Investigação científica - Investigar e analisar \textbf{problemas} - Levantar e testar hipóteses - Avaliar e prever efeitos de intervenções nos ecossistemas, e seus impactos nos seres vivos e no corpo humano, com base nos mecanismos de manutenção da vida, nos ciclos da matéria e nas transformações e transferências de energia, utilizando representações e \textbf{simulações} sobre tais fatores, com ou sem o uso de dispositivos e \textbf{aplicativos} \textbf{digitais} ( como softwares de \textbf{simulação} e de realidade \textbf{virtual}, entre outros ). - Participar de pesquisas científicas que investiguem \textbf{problemas} da cidade relacionados aos impactos da ação humana na natureza ( e ciclos ) que envolvam solo, ar e água.
Objetos do conhecimento - Equilíbrio ambiental - Impactos ambientais na fauna e flora relacionados à poluição e à contaminação do ar, da água e do solo - Consequências da urbanização - Sucessão ecológica - Atmosfera terrestre e alterações climáticas - Ciclos biogeoquímicos ( água, carbono, nitrogênio, oxigênio e fósforo ) - Fenômenos atmosféricos naturais e causados pela ação humana ( vento, umidade, precipitação, formação de nuvens, ciclones, inversão térmica, chuvas ácidas, desastres naturais ) - Poluição atmosférica local e global ( causas e principais impactos \textbf{socioambientais} ), alterações atmosféricas e efeitos na saúde humana ( doenças respiratórias, mutações genéticas, câncer ) - Calorimetria e hidrostática ( umidade do ar ; \textbf{capacidade} térmica, calor específico, temperatura, pressão, densidade, processos de propagação de calor ) aplicadas à climatologia - Termoquímica ; reações fotoquímicas - Efeito estufa e aquecimento global - Negacionismo científico em relação às mudanças climáticas - Recursos hídricos e cuidados com o solo - Ciclos biogeoquímicos ( água, carbono, nitrogênio, oxigênio e fósforo ) - Propriedades e estados físicos da água ( em nível de estrutura molecular ) - Reservas hídricas locais, utilização e condições de preservação - causas antrópicas da poluição de aquíferos ( áreas urbanas sem redes de esgoto ; atividades industriais ; resíduos sólidos ; atividades agrícolas ; extrativismo mineral ; acidentes ambientais e tanques enterrados ) - Legislação ambiental – plano nacional dos recursos hídricos - Solo : utilização, cuidados, manejo adequado ( controle biológico de pragas, drenagem, agricultura orgânica... ) - Técnicas de recuperação de áreas degradadas - Fontes de poluição do solo e da água e suas consequências ( extração de minérios, acúmulo de metais pesados, uso de agrotóxicos, desmatamentos, lançamento de lixo, materiais radioativos, lixo tecnológico ) - Cadeias alimentares e bioacumulação - Enchentes e alagamentos – causas e consequências em nível local Fonte : Elaboração dos \textbf{autores} ( 2020 ).
Unidade curricular III : Saneamento básico e gerenciamento de resíduos Objetivo : Promover a investigação acerca do saneamento básico e do gerenciamento dos resíduos nas cidades, com base na Constituição, bem como reconhecer a importância das ações individuais e coletivas, para a melhoria das condições ambientais, sociais, de saúde e de qualidade de vida da população local a fim de propor alternativas que melhorem a prestação desses serviços em nível local.
Carga horária : Total de 40 h ou 60 h, a \textbf{depender} da \textbf{matriz} curricular em funcionamento na Unidade Escolar. \textbf{Quadro} 28 - Saneamento básico e gerenciamento de resíduos Eixo estruturante/Palavras-chave das habilidades específicas - Empreendedorismo - Avaliar possíveis soluções - Desenvolver projetos - Criar propostas articuladas ao projeto de vida - Avaliar os benefícios e os riscos à saúde e ao ambiente, considerando a composição, a toxicidade e a reatividade de diferentes materiais e produtos, como também o nível de \textbf{exposição} a eles, posicionando -se criticamente, propondo soluções individuais e/ou coletivas para seus usos e descartes responsáveis. - Investigar e analisar os efeitos de programas de infraestrutura e demais serviços básicos ( saneamento, energia elétrica, transporte, telecomunicações, cobertura vacinal, atendimento primário à saúde e produção de alimentos, entre outros ) ; identificar necessidades locais e/ou regionais em relação a esses serviços a fim de avaliar e/ou promover ações que contribuam para a melhoria da qualidade de vida e das condições de saúde da população.
Objetos do conhecimento - Fluidodinâmica aplicada ao estudo - Escoamento de fluidos, equação da continuidade, equação de Bernoulli, viscosidade de fluidos, modelos hidrodinâmicos - Tratamento da água - Captação : fontes locais de abastecimento de água - Estações e etapas de tratamento de água ( processos físicos, biológicos e químicos envolvidos ) - Distribuição, custo e consumo de água - Alternativas ao tratamento da água onde não há estações de tratamento da água - Tratamento de esgoto - Importância do tratamento de esgoto do ponto de \textbf{vista} ambiental e da saúde pública. - Principais doenças associadas à água contaminada e ao contato com esgotos em nível local - Alternativas para o tratamento de águas residuárias - Estações de tratamento de esgotos - Etapas de tratamento do esgoto ( processos biológicos, químicos e físico-químicos envolvidos ) e disposição \textbf{final} de águas residuárias e lodos gerados - Gerenciamento de resíduos - Principais características dos resíduos sólidos ( inflamabilidade, corrosividade, reatividade, toxicidade, patogenicidade, radioatividade, biodegradabilidade, combustibilidade ) - Separação e destinação dos resíduos em nível local - Programas locais para transporte, coleta, redução e reciclagem - Tratamento e disposição \textbf{final} dos resíduos ( compostagem, incineração, pirólise, aterros ) - Possibilidades de aproveitamento de resíduos orgânicos para a produção de energia ( biodigestores ) Fonte : Elaboração dos \textbf{autores} ( 2020 ).
Unidade curricular IV : Planejamento ambiental das cidades - planos diretores, mobilidade urbana e sustentabilidade Objetivo : Oferecer \textbf{subsídios} para a \textbf{análise} de questões que interferem diretamente na dinâmica das cidades, bem como instrumentos básicos de política de desenvolvimento e expansão urbana.
Através da educação ambiental, promover a reflexão acerca de atitudes, valores, inovações, \textbf{tecnologias} e formas de construir uma sociedade mais sustentável.
Carga horária : Total de 40 h ou 60 h, a \textbf{depender} da \textbf{matriz} curricular em funcionamento na Unidade Escolar. \textbf{Quadro} 29 - Planejamento ambiental das cidades - planos diretores, mobilidade urbana e sustentabilidade Eixo estruturante/Palavras-chave das habilidades específicas - Mediação e intervenção sociocultural - Identificar e explicar fenômenos. - Selecionar e \textbf{mobilizar} conhecimentos científicos - Propor e testar intervenções socioculturais e ambientais Habilidades de \textbf{aprofundamento} da área - Discutir a importância da preservação e conservação da biodiversidade, considerando \textbf{parâmetros} \textbf{qualitativos} e quantitativos, e avaliar os efeitos da ação humana e das políticas ambientais para a garantia da sustentabilidade do planeta. - Construir questões, elaborar hipóteses, previsões e estimativas, empregar instrumentos de medição e representar e interpretar modelos explicativos, dados e/ou resultados experimentais para construir, avaliar e justificar conclusões no enfrentamento de situações-problema sob uma perspectiva científica. - Analisar o plano diretor municipal, identificando as possibilidades de participação da comunidade nas decisões e relacionar com as mudanças implementadas a partir do plano. - Identificar, do ponto de \textbf{vista} ambiental, as profissões relacionadas à dinâmica das cidades.
Objetos do conhecimento - Estatuto da cidade - Plano diretor municipal : \textbf{atores} envolvidos, o que está contemplado, mudanças e motivações, proteção ao meio ambiente, cumprimento do plano em sua \textbf{totalidade} - Poluição sonora - Profissões relacionadas às questões ambientais - Transporte e mobilidade urbana : - Leis da termodinâmica aplicadas aos meios de transporte - Principais meios de transporte utilizados - públicos e coletivos ( carros movidos a gasolina, diesel, álcool, balsas, motocicletas, carros elétricos, bicicletas, entre outros ) - Uso de combustíveis fósseis X fontes alternativas e renováveis ( biocombustíveis ) de energia - Mecânica clássica ( Leis de Newton, conservação da energia, conservação da quantidade de movimento ) e segurança no trânsito : cinto de segurança, tipos e funcionamento de freios, importância da manutenção de um veículo, acidentes de trânsito, legislação e impactos na saúde pública - Mobilidade urbana : plano, \textbf{problemas} locais, equipamentos de controle ( radares e semáforos ), acessibilidade, “ encurtamento ” de distâncias, química verde e sustentabilidade - Educação ambiental e Sustentabilidade : - Lei de proteção aos animais - Princípios da Sustentabilidade - Legislação ambiental e agricultura ( código florestal, áreas de preservação permanente, reserva legal, unidades de conservação, produção agroecológica, etc. ) - Política dos R’s - Obsolescência programada e consumo consciente - Fake News ( notícias falsas ) relacionadas ao meio ambiente - Casas eficientes e sustentáveis : iluminação, conforto térmico, geração de energia elétrica, distribuição dos cômodos, economia de recursos, matérias-primas para as construções e as mobílias ( tipos, reaproveitamento, racionalização ) - Cidades e áreas verdes Fonte : Elaboração dos \textbf{autores} ( 2020 ). 3.2.4 Orientações \textbf{Metodológicas} Uma \textbf{abordagem} que se alinha com a proposta desta \textbf{trilha} é a \textbf{abordagem} \textbf{Ciência}, \textbf{Tecnologia} e Sociedade ( CTS ), que \textbf{implica} a utilização de \textbf{eventos}, temas, \textbf{tecnologias} e fenômenos para, a partir deles, tratar dos conceitos.
No caso da \textbf{trilha} aqui proposta, o tema são as cidades e os conceitos específicos da área utilizados para se compreender as dinâmicas mais amplas.
A \textbf{abordagem} CTS possibilita desenvolver nos(as) estudantes a \textbf{capacidade} de problematizar questões de \textbf{ciência} e \textbf{tecnologia}, produzir conhecimentos, tomar decisões \textbf{argumentadas} em \textbf{fundamentos} científicos, com autonomia para exercer a cidadania consciente.
Entre os \textbf{argumentos} em defesa da educação em CTS estão a superação do ensino dogmático, a \textbf{compreensão} das relações CTS, a tomada de decisões, em nível pessoal e coletivo, fundamentadas nas consequências da atuação do ser humano no meio e a busca de sentido para a \textbf{ciência} da escola ( STRIEDER et al., 2017 ).
A \textbf{abordagem} CTS, em nível de racionalidade científica, admite o conhecimento científico como uma produção para compreender o mundo, a importância de conhecer conceitos e processos da \textbf{ciência} para tomar decisões sobre \textbf{problemas} relacionados a ela, as implicações dos conhecimentos científicos na sociedade e discussões sobre a natureza da \textbf{Ciência}.
No nível do desenvolvimento tecnológico, relaciona funcionamento, produção, processos, uso e adequação social de \textbf{tecnologias}.
A participação social considera \textbf{problemas} gerados pela \textbf{Ciência} e \textbf{Tecnologia}, a intervenção social, a educação ambiental, a tomada de decisões sobre o uso de produtos ou \textbf{tecnologias} e a participação social no âmbito da \textbf{Ciência} e \textbf{Tecnologia}.
Diversas são as metodologias que se podem utilizar no desenvolvimento da \textbf{trilha}, a exemplo das citadas no texto da formação geral básica.
Dentre elas está o desenvolvimento de projetos interdisciplinares ou de investigação científica que podem ser elaborados pelos professores da área, ou com os(as) estudantes, a partir de situações ou \textbf{problemas} do seu entorno, como, por exemplo, tratamento de esgotos, enchentes, alagamentos, mobilidade urbana, entre outros.
Os projetos podem ser uma forma de estabelecer um elo entre os conceitos científicos e a sociedade, ligando estes conceitos a determinado contexto da realidade dos(as) estudantes e não ao contexto dos componentes curriculares.
É importante o desenvolvimento de todas as etapas do projeto, que incluem elaboração, desenvolvimento, avaliação e comunicação de resultados.
A comunicação pode acontecer em sala de \textbf{aula}, na escola ou na comunidade, podendo converter -se em soluções para \textbf{problemas} locais e transformar a realidade, no sentido de intervenção social.
Uma proposta para o desenvolvimento de projetos interdisciplinares encontra -se nas ‘ ilhas interdisciplinares de racionalidade ’ ( NEHRING et al., 2002 ), que envolvem a elaboração de um projeto, dividido em oito etapas, a partir de um contexto, ou \textbf{problema}.
O projeto prevê a elaboração de um produto destinado a um público, escolhido por quem elabora o projeto.
É importante que as metodologias de ensino sejam diversificadas, porém, pautadas em experiências de aprendizagem que provoquem os(as) estudantes, que os(as) motivem a se tornar protagonistas no processo vivido, sendo mais criativos(as) e empreendedores(as ). \textbf{Convém} \textbf{lembrar} que \textbf{vivemos} numa sociedade de informação, que dispõe de recursos \textbf{digitais} que podem contribuir no engajamento dos(as) estudantes e valorizar os diferentes estilos e valores comunicacionais. \textbf{Sugerimos}, a seguir, algumas atividades que podem \textbf{auxiliar} os professores no desenvolvimento das unidades curriculares, como : a ) desenvolvimento de projetos, experimentos e protótipos ; b ) levantamento dos principais aparelhos eletrônicos e eletrodomésticos existentes nas residências dos(as) estudantes, e respectivas transformações de energia, cálculo de consumo de energia elétrica ; c ) explanação, \textbf{debate} e utilização de materiais audiovisuais ( \textbf{vídeos}, documentários, \textbf{entrevistas}, podcast, entre outros ) ; d ) pesquisas, \textbf{entrevistas} e \textbf{palestras} com profissionais relacionadas a questões ambientais ; e ) produção de materiais de divulgação científica ; f ) visitas a aterros sanitários e lixões ; g ) \textbf{análise} de artigos científicos relacionados com a temática e \textbf{sugeridos} pelos professores da área.
AVALIAÇÃO A avaliação será contínua, processual e formativa, com prevalência dos aspectos \textbf{qualitativos} sobre os quantitativos, propiciando movimentos de reflexão sobre o próprio processo de ensino e aprendizagem, com \textbf{análise} das dificuldades, das lacunas e na busca constante da efetivação do conhecimento.
Nesse cenário, em que o projeto de vida do(a) estudante está em construção, a avaliação deve ser pautada no processo que ocorre, nas habilidades desenvolvidas e nos objetivos e metas a serem alcançados, no comprometimento com os planos de ação idealizados e na recuperação paralela.

% matched lemmas: abordagem, análise, aplicativo, aprofundamento, aproximar, argumentar, argumento, articulado, ator, atual, aula, autor, auxiliar, caminho, capacidade, ciência, compreensão, contemporâneo, convir, crítico, debate, depender, desempenhar, dialógico, digital, divulgar, econômico, entrevista, escopo, evento, exigir, exposição, figurar, final, foco, fundamento, horizonte, humanidade, implicar, influenciar, influência, inúmero, lado, lembrar, maneira, matriz, metodológico, mobilizar, palestra, parâmetro, potencial, preocupar, primordial, problema, quadro, qualitativo, reflexivo, segar, sequência, simulação, socioambiental, subsídio, sugerir, supor, tanto, tarefa, tecnologia, totalidade, transformador, trilha, virtual, vista, vivar, vídeo
\end{textsample}
